% Level 2, Week 1 -- homework
\homeworkweek{1}{Welcome Back, Python}{15--20 minutes. Do Core first; Challenge and Creative are optional.}
\namedateline

\begin{core}[Type detective (paper, 3 minutes)]
Write the type of each value: \textbf{int}, \textbf{float}, \textbf{str} or \textbf{bool}.
\begin{center}
\renewcommand{\arraystretch}{1.6}
\begin{tabular}{|l|p{2.2cm}||l|p{2.2cm}|}
\hline
\texttt{42} & & \texttt{"42"} & \\ \hline
\texttt{4.2} & & \texttt{False} & \\ \hline
\texttt{"False"} & & \texttt{-8} & \\ \hline
\end{tabular}
\end{center}
\end{core}

\begin{core}[Dragon Door (Thonny, 10 minutes)]
Open \texttt{excercises/week01/hw-dragon-door/starter.py}. The dragon door opens if:
\begin{itemize}[leftmargin=1.2em]
  \item you have the golden key \textbf{and} you are level 3 or more, \textbf{or}
  \item you are level 10 or more (legends do not need keys).
\end{itemize}
\begin{enumerate}[leftmargin=1.4em]
  \item TODO 1: change the \texttt{if} line in \texttt{door\_opens()}. Brackets help: \texttt{(has\_key and level >= 3) or level >= 10}
  \item TODO 2: make the welcome message an f-string that says the hero name.
  \item Run it three times and tick what happened:\par
  \begin{tabular}{@{}lll@{}}
  key yes, level 5: & opens \tickbox & shut \tickbox \\
  key yes, level 2: & opens \tickbox & shut \tickbox \\
  key no, level 12: & opens \tickbox & shut \tickbox \\
  \end{tabular}
  \item Open \texttt{check\_my\_work.py} and press \menu{Run}.
\end{enumerate}
\end{core}

\begin{challenge}[Predict, then check]
Write what each line prints. Then type it into the Thonny Shell to check.
\begin{pycode}
print(3 > 2 and 5 > 9)
print(3 > 2 or 5 > 9)
print(not 3 > 2)
print("7" + "7")
print(int("7") + 7)
\end{pycode}
\answerlines{3}
\end{challenge}

\begin{creative}[Design a hero]
Invent a hero for a game (a made-up name only). Draw them in the box and write a hero card with at least
one value of each type: a str, an int, a float and a bool. Next week you can type it into the Hero Card program.
\answerspace{4cm}
\end{creative}

\begin{selfcheck}
  \item I can name the four types and give an example of each.
  \item I know that \texttt{input()} gives back text, and how \texttt{int()} fixes that.
  \item My Dragon Door opens for the right heroes and stays shut for the others.
  \item I used a made-up hero name, not my real name.
\end{selfcheck}
